\section{Threats to Validity}

\myparagraph{Construct and internal validity}
Our verifiability rates are lower bounds. We count an artifact as unverified when \toolname cannot recover its rebuild metadata or reproduce it with canonical build commands. Thus, the rates reflect both ecosystem verifiability and \toolname's verifier model, and environment reconstruction; a more complete verifier may verify more artifacts.

Our L3 results depend on a manually curated catalog of non-executable differences, built from ecosystem documentation and analysis of non-L1 cases. This catalog may be too restrictive or permissive, so we interpret L3 verification as evidence of source-to-artifact correspondence, not as proof that all normalized differences are security-neutral.

Our RQ2 root-cause analysis may also misattribute failures. We used an automated agent to scale manual coding and accepted hypothesized causes when corrective changes improved equivalence. Such changes may identify contributing causes rather than the only or true cause. Moreover, the taxonomy is derived from provenance-bearing artifacts, whose failures may differ from non-attested artifacts.

\myparagraph{External validity and reliability}
Our study evaluates only the most recent artifact of each package~(\cref{subsubsec:setup-artifact-selection}), so it does not capture version-to-version variation or temporal drift in tools, dependencies, and attestation adoption.

We restrict rebuilds to packages with recoverable public source repositories, excluding 14.7-42.4\%~\cite{PypiNpmGHAccessibility, tsakpinis2026investigatingnotablemetadatapractices}. Reported rates are therefore conditional on source recovery and may overestimate ecosystem-wide verifiability. Because we oversample attested packages, ecosystem rates are stratum-weighted rather than population-weighted, and we may miss packages with attestations outside registry-provided pathways.

Finally, our OSS-Rebuild comparison reflects one tool version and our operation of it, which may miss capabilities requiring additional configuration or expertise.



% Because we evaluate only the latest artifact for each package (\cref{subsubsec:setup-artifact-selection}), our results may not capture version-to-version variation in verifiability.
% \comment[Oreofe]{we can state that if we did a horizontal study (i.e., looking at multiple versions of packages), it might not provide us much detail.  Maybe if we had studied the verifiability of packages improved over time, (if there time we can take a small sample (popular packages) and do it?  }

% We may also miss packages that create and store attestations manually and not through the registry-provided pathway.